\begin{table}[t]
\centering
\renewcommand{\arraystretch}{1.2}
\caption{Comparison with other qLDPC decoders.}
\label{tab:intro_comparison}
\resizebox{\linewidth}{!}{
\begin{minipage}{\linewidth}
\begin{tabular}{|l|c|c|c|c|}
\toprule[1pt]
    \textbf{\makecell{Decoder Type}} & \textbf{\makecell{Decoding\\Accuracy}} & \textbf{\makecell{Time\\Complexity}} & \textbf{\makecell{Technical\\Feature}} & \textbf{Principle} \\
    \hline
    AFS \cite{das2022afs} & \bad & \good & \makecell{Union-find} & MWD \\
    \hline
    Astrea \cite{vittal2023astrea} & \good & \medium & \makecell{Graph\\Matching}  & MWD \\
    \hline
    BP \cite{bp} & \bad & \good & \makecell{Message\\Passing}  & MWD \\
    \hline
    BP+OSD \cite{bposd} & \good & \bad & \makecell{Gaussian\\Elimination}  & MWD \\
    \hline
    BP+LSD \cite{bplsd} & \medium & \medium & \makecell{Matrix\\Factorization}  & MWD \\
    \hline
    \textbf{\papername} & \good & \good & \makecell{Tensor\\Compression} & \textbf{MLD} \\
\bottomrule[1pt]
\end{tabular}

\vspace{0.15cm}
\raggedright
\footnotesize
$^*$MWD refers to Minimum Weight Decoding, while MLD means Maximum Likelihood Decoding. \\
$^*$The symbol $\odot$ indicates a medium level.
\end{minipage}
}
\vspace{-0.4cm}
\end{table}

